\chapter{Graph Theory — Karger Min Cut}
%%% generated by the blueprint pipeline from TCSlib.GraphTheory.KargerMinCut
%%%%%%%%%%%%%%%%%%%%%%%%%%%%%%%%%%%%%%%%%%%%%%%%%%%%%%%%%%%%%%%%%%%%%%%%%%%%%%%
\section{Overview}

This file develops a multigraph model (a \texttt{Multiset} of unordered pairs, since
Mathlib's \texttt{SimpleGraph} has no parallel edges) together with edge contraction,
cuts and their sizes, and the structural lemmas behind Karger's randomized contraction
algorithm: contraction of a non-crossing edge preserves a minimum cut, every vertex
degree bounds the min-cut size, the handshake lemma $\sum_v \deg(v) = 2\abs{E}$, the
telescoping survival product $2/(n(n-1))$, and the standard corollaries on repetition
and on the number of (approximate) minimum cuts.

%%%%%%%%%%%%%%%%%%%%%%%%%%%%%%%%%%%%%%%%%%%%%%%%%%%%%%%%%%%%%%%%%%%%%%%%%%%%%%%
\section{Declarations}

\begin{definition}[Multigraph]
\lean{Multigraph}
\label{Multigraph}
\leanok
A multigraph on a vertex type $\alpha$ consists of a finite vertex set together with a
multiset of unordered pairs (elements of $\mathrm{Sym}_2\,\alpha$) as edges, subject to
two conditions: every endpoint of every edge is a live vertex, and no edge is a
self-loop (edges are loopless, since loops are deleted after each contraction).
\end{definition}

\begin{definition}[Edge count of a multigraph]
\lean{Multigraph.edgeCount}
\label{Multigraph.edgeCount}
\leanok
\uses{Multigraph}
The number of edges of $G$ counted with multiplicity, i.e.\ the cardinality of the
edge multiset.
\end{definition}

\begin{definition}[Vertex count of a multigraph]
\lean{Multigraph.vertexCount}
\label{Multigraph.vertexCount}
\leanok
\uses{Multigraph}
The number of vertices of $G$, i.e.\ the cardinality of its vertex \texttt{Finset}.
\end{definition}

\begin{definition}[Contraction of an edge]
\lean{Multigraph.contract}
\label{Multigraph.contract}
\leanok
\uses{Multigraph}
Given an edge $e \in E(G)$ with endpoints $u$ and $v$, the contracted multigraph
$G/e$ is obtained by replacing every occurrence of $v$ by $u$ in all edges, deleting
the resulting self-loops (in particular all parallel copies of $e$), and removing $v$
from the vertex set. Parallel edges between the merged vertex and the rest of the graph
are retained with multiplicity.
\end{definition}

\begin{lemma}[Contraction reduces the vertex count by one]
\lean{contract_vertex_count}
\label{contract_vertex_count}
\leanok
\uses{Multigraph, Multigraph.contract, Multigraph.vertexCount}
\difficulty{2}
Let $G$ be a multigraph and let $e$ be an edge of $G$ that is not a self-loop, so that
its two endpoints are distinct. Writing $G/e$ for the multigraph obtained by contracting
$e$, its vertex count satisfies
\[
  \abs{V(G/e)} = \abs{V(G)} - 1,
\]
where the subtraction is truncated at $0$.
\end{lemma}

\begin{lemma}[Contraction removes at least one edge]
\lean{contract_edge_count_le}
\label{contract_edge_count_le}
\leanok
\uses{Multigraph, Multigraph.contract, Multigraph.edgeCount}
\difficulty{5}
Let $G$ be a multigraph and let $e \in E(G)$ be one of its edges. Then the multigraph
$G/e$ obtained by contracting $e$ satisfies
\[
  \abs{E(G/e)} \le \abs{E(G)} - 1,
\]
where both edge counts are taken with multiplicity; in particular, contracting an edge
never increases the number of edges and eliminates at least the contracted edge itself.
\end{lemma}

\begin{definition}[Cut of a multigraph]
\lean{MulCut}
\label{MulCut}
\leanok
\uses{Multigraph}
A cut of $G$ is a subset $S \subseteq V(G)$ that is nonempty and whose complement
$V(G) \setminus S$ inside $V(G)$ is also nonempty; the two nonemptiness conditions rule
out the trivial partitions.
\end{definition}

\begin{definition}[Crossing edges of a cut]
\lean{MulCut.crossingEdges}
\label{MulCut.crossingEdges}
\leanok
\uses{MulCut, Multigraph}
The multiset of edges of $G$ that cross the cut $C$, i.e.\ those edges having exactly
one endpoint in $S$, counted with multiplicity. It is obtained by filtering the edge
multiset by the condition that exactly one of the two representative endpoints lies
in $S$.
\end{definition}

\begin{definition}[Crossing predicate for an unordered edge]
\lean{Crosses}
\label{Crosses}
\leanok
For a finite set $S$ and an unordered pair $e$, the predicate $\texttt{Crosses}\;S\;e$
holds when some element of $e$ lies in $S$ and some element of $e$ lies outside $S$.
This is a membership-based formulation independent of a choice of representative.
\end{definition}

\begin{lemma}[Crossing via a chosen representative of the edge]
\lean{crosses_iff_out}
\label{crosses_iff_out}
\leanok
\uses{Crosses}
\difficulty{4}
Let $S$ be a finite subset of a type $\alpha$, and let $e$ be an unordered pair of
elements of $\alpha$; write $(x,y)$ for a chosen ordered representative of $e$. Then $e$
crosses $S$ — that is, some element of $e$ lies in $S$ and some element of $e$ lies
outside $S$ — if and only if either $x \in S$ and $y \notin S$, or $x \notin S$ and $y
\in S$.
\end{lemma}

\begin{lemma}[Crossing edges as the membership-based edge filter]
\lean{crossingEdges_eq_filter_crosses}
\label{crossingEdges_eq_filter_crosses}
\leanok
\uses{Crosses, MulCut, MulCut.crossingEdges, Multigraph, crosses_iff_out}
\difficulty{2}
Let $G$ be a multigraph and let $C$ be a cut of $G$ determined by a vertex subset $S$.
Then the multiset of edges crossing $C$ equals the sub-multiset of edges $e$ of $G$ for
which some endpoint of $e$ lies in $S$ while some endpoint of $e$ lies outside $S$,
counted with multiplicity.
\end{lemma}

\begin{definition}[Size of a cut]
\lean{MulCut.size}
\label{MulCut.size}
\leanok
\uses{MulCut, MulCut.crossingEdges, Multigraph}
The size of a cut $C$ is the number of crossing edges, counted with multiplicity.
\end{definition}

\begin{definition}[Minimum cut]
\lean{IsMulMinCut}
\label{IsMulMinCut}
\leanok
\uses{MulCut, MulCut.size, Multigraph}
A cut $C$ of $G$ is a minimum cut when its size is at most the size of every cut $C'$
of $G$.
\end{definition}

\begin{definition}[Cut induced on a contracted multigraph]
\lean{MulCut.contractedCut}
\label{MulCut.contractedCut}
\leanok
\uses{MulCut, MulCut.crossingEdges, Multigraph, Multigraph.contract}
If $e \in E(G)$ does not cross the cut $C$, then $C$ induces a cut of the contracted
multigraph $G/e$: the vertex side is the image of $S$ after merging the endpoints of
$e$, and both nonemptiness conditions continue to hold because $e$ has both endpoints
on the same side.
\end{definition}

\begin{lemma}[Contracting a non-crossing edge preserves cut size]
\lean{cut_size_preserved}
\label{cut_size_preserved}
\leanok
\uses{Crosses, MulCut, MulCut.contractedCut, MulCut.crossingEdges, MulCut.size, Multigraph, Multigraph.contract, crosses_iff_out, crossingEdges_eq_filter_crosses}
\difficulty{6}
Let $G$ be a multigraph with vertex set $V(G)$ and edge multiset $E(G)$, and let $C$ be
a cut of $G$, given by a nonempty proper subset $S \subseteq V(G)$ with nonempty
complement. Suppose $e \in E(G)$ is an edge that does not cross $C$, i.e.\ both
endpoints of $e$ lie on the same side of the partition. Then the cut induced on the
contracted multigraph $G/e$ by merging the endpoints of $e$ has the same size as $C$:
the number of edges crossing the induced cut in $G/e$, counted with multiplicity, equals
the number of edges crossing $C$ in $G$.
\end{lemma}

\begin{lemma}[Non-crossing contraction preserves minimality]
\lean{mincut_preserved_of_non_crossing}
\label{mincut_preserved_of_non_crossing}
\leanok
\uses{IsMulMinCut, MulCut, MulCut.contractedCut, MulCut.crossingEdges, MulCut.size, Multigraph, Multigraph.contract, count_bind_endpoints_eq_card_filter, crossingEdges_eq_filter_crosses, cut_size_preserved}
\difficulty{6}
Let $G$ be a multigraph, and let $C$ be a minimum cut of $G$, given by a vertex subset
$S$ that is nonempty and has nonempty complement in $V(G)$. Suppose $e \in E(G)$ is an
edge that does not cross $C$, so that both endpoints of $e$ lie on the same side of the
partition. Then the cut induced by $C$ on the contracted multigraph $G/e$ — obtained by
merging the endpoints of $e$ and intersecting $S$ with the surviving vertices — is again
a minimum cut, of $G/e$.
\end{lemma}

\begin{definition}[One contraction step of Karger's algorithm]
\lean{kargerStep}
\label{kargerStep}
\leanok
\uses{Multigraph, Multigraph.contract, Multigraph.edgeCount}
Given a multigraph $G$ with at least one edge, a step of the algorithm is modelled as a
function from an index $i \in \mathrm{Fin}\,\abs{E(G)}$ — thought of as a uniformly
random edge choice — to the multigraph obtained by contracting the $i$-th edge of $G$.
\end{definition}

\begin{definition}[A full contraction run]
\lean{kargerRun}
\label{kargerRun}
\leanok
\uses{Multigraph, Multigraph.edgeCount, Multigraph.vertexCount}
A run of the algorithm on $G_0$ is indexed by a sequence of edge choices, one for each
of the $\abs{V(G_0)} - 2$ contraction steps, and produces a multigraph. The current
Lean definition is a placeholder that returns $G_0$; the state-threading version is not
yet formalized.
\end{definition}

\begin{definition}[Termination condition of the algorithm]
\lean{KargerOutput}
\label{KargerOutput}
\leanok
\uses{Multigraph, Multigraph.vertexCount}
The predicate stating that $G$ is a valid output of the algorithm, namely that it has
exactly two vertices; its remaining edges are then exactly the edges of the output cut.
\end{definition}

\begin{definition}[Degree of a vertex]
\lean{Multigraph.degree}
\label{Multigraph.degree}
\leanok
\uses{Multigraph}
The degree of $v$ in a loopless multigraph $G$ is the number of edges incident to $v$,
counted with multiplicity.
\end{definition}

\begin{lemma}[Membership in an unordered pair meets a chosen representative]
\lean{eq_out_or_eq_out_of_mem}
\label{eq_out_or_eq_out_of_mem}
\leanok
\difficulty{2}
Let $e$ be an unordered pair of elements of a type $\alpha$ with decidable equality, and
let $(a, b)$ be the ordered pair chosen as a representative of $e$. If an element $z$
belongs to $e$, then $z = a$ or $z = b$.
\end{lemma}

\begin{lemma}[Endpoint multiplicity equals incidence count]
\lean{count_bind_endpoints_eq_card_filter}
\label{count_bind_endpoints_eq_card_filter}
\leanok
\difficulty{4}
Let $m$ be a multiset of unordered pairs (edges) drawn from a type $\alpha$, and fix an
element $v \in \alpha$. For each edge $e$, let $\{e\}_\ast$ denote the finite set of its
endpoints — a singleton when $e$ is a loop and a two-element set otherwise — and form
the multiset obtained by concatenating $\{e\}_\ast$ over all $e$ in $m$, counting
multiplicities of $m$. Then the multiplicity of $v$ in this concatenated multiset equals
the number of edges $e$ of $m$ for which $v$ is an endpoint of $e$.
\end{lemma}

\begin{lemma}[Size of a singleton cut equals the vertex degree]
\lean{singleton_cut_size_eq_degree}
\label{singleton_cut_size_eq_degree}
\leanok
\uses{MulCut, MulCut.crossingEdges, MulCut.size, Multigraph, Multigraph.degree, crosses_iff_out, crossingEdges_eq_filter_crosses, eq_out_or_eq_out_of_mem}
\difficulty{5}
Let $G$ be a loopless multigraph with vertex set $V(G)$, and let $v \in V(G)$ be a
vertex whose deletion leaves at least one other vertex, i.e. $V(G) \setminus \{v\}$ is
nonempty. Then the singleton $S = \{v\}$ is a cut of $G$, and the size of this cut — the
number of edges with exactly one endpoint in $\{v\}$, counted with multiplicity — equals
the degree $\deg(v)$ of $v$.
\end{lemma}

\begin{lemma}[Deleting a vertex from a cuttable multigraph]
\lean{erase_nonempty_of_cut}
\label{erase_nonempty_of_cut}
\leanok
\uses{MulCut, Multigraph}
\difficulty{3}
Let $G$ be a multigraph that admits a cut. Then for every vertex $v$, the vertex set of
$G$ with $v$ removed, $V(G) \setminus \{v\}$, is nonempty.
\end{lemma}

\begin{lemma}[Minimum cut is at most every vertex degree]
\lean{mincut_le_degree}
\label{mincut_le_degree}
\leanok
\uses{IsMulMinCut, MulCut, MulCut.size, Multigraph, Multigraph.degree, erase_nonempty_of_cut, singleton_cut_size_eq_degree}
\difficulty{3}
Let $G$ be a loopless multigraph, let $C$ be a minimum cut of $G$, and let $v$ be any
vertex of $G$. Then the size of $C$ is at most the degree of $v$, that is,
$\mathrm{size}(C) \le \deg(v)$.
\end{lemma}

\begin{lemma}[Endpoints of edges are vertices]
\lean{endpoints_mem_vertices}
\label{endpoints_mem_vertices}
\leanok
\uses{Multigraph}
\difficulty{2}
Let $G$ be a multigraph on a vertex type $\alpha$, with vertex set $V(G)$ and edge
multiset $E(G)$. If a point $x \in \alpha$ occurs as an endpoint of some edge of $G$ —
that is, $x$ lies in one of the unordered pairs making up $E(G)$ — then $x \in V(G)$.
\end{lemma}

\begin{lemma}[Summing multiplicities over a covering set]
\lean{sum_counts_eq_card_of_subset}
\label{sum_counts_eq_card_of_subset}
\leanok
\difficulty{5}
Let $m$ be a finite multiset with elements drawn from a type $\alpha$, and let $s$ be a
finite set of elements of $\alpha$ that contains every element occurring in $m$. Writing
$\mathrm{count}_m(a)$ for the multiplicity of $a$ in $m$, the multiplicities over $s$
sum to the total size of the multiset:
\[
  \sum_{a \in s} \mathrm{count}_m(a) = \abs{m}.
\]
\end{lemma}

\begin{lemma}[Handshake lemma for multigraphs]
\lean{sum_degrees_eq_twice_edgeCount}
\label{sum_degrees_eq_twice_edgeCount}
\leanok
\uses{Multigraph, Multigraph.degree, Multigraph.edgeCount, count_bind_endpoints_eq_card_filter, endpoints_mem_vertices, sum_counts_eq_card_of_subset}
\difficulty{4}
Let $G$ be a loopless multigraph, given by a finite vertex set $V$ together with a
multiset $E$ of edges, each edge an unordered pair of distinct vertices of $V$. For a
vertex $v$, let $\deg(v)$ be the number of edges of $G$ incident to $v$, counted with
multiplicity. Then
\[
  \sum_{v \in V} \deg(v) = 2\,\abs{E},
\]
where $\abs{E}$ is the number of edges counted with multiplicity.
\end{lemma}

\begin{lemma}[Edge count lower bound from a minimum degree]
\lean{edge_count_lower_bound}
\label{edge_count_lower_bound}
\leanok
\uses{Multigraph, Multigraph.degree, Multigraph.edgeCount, Multigraph.vertexCount, sum_degrees_eq_twice_edgeCount}
\difficulty{3}
Let $G$ be a loopless multigraph with finite vertex set $V$ and edge multiset $E$, and
let $k$ be a natural number. If every vertex $v \in V$ has degree at least $k$ in $G$,
then
\[
  k \cdot \abs{V} \le 2\,\abs{E},
\]
where $\abs{V}$ is the number of vertices and $\abs{E}$ is the number of edges counted
with multiplicity.
\end{lemma}

\begin{definition}[Survival probability of one step]
\lean{survivalProb}
\label{survivalProb}
\leanok
For a cut size $c$ and a current edge count $m$, the probability that a uniformly chosen
edge is not one of the $c$ cut edges, namely
\[
  1 - \frac{c}{m}.
\]
\end{definition}

\begin{lemma}[Edge-count lower bound at each contraction step]
\lean{edge_count_invariant}
\label{edge_count_invariant}
\leanok
\uses{IsMulMinCut, MulCut, MulCut.size, Multigraph, Multigraph.edgeCount, Multigraph.vertexCount, edge_count_lower_bound, mincut_le_degree}
\difficulty{3}
Let $G_0$ be a multigraph, and let $C_0$ be a minimum cut of $G_0$, so that the number
of edges crossing $C_0$ (counted with multiplicity), its size, is no larger than that of
any other cut. Then for every natural number $i$ with $i < \abs{V(G_0)} - 1$,
\[
  \mathrm{size}(C_0)\cdot\bigl(\abs{V(G_0)} - i\bigr) \;\le\; 2\,\abs{E(G_0)},
\]
where $\abs{V(G_0)}$ is the vertex count and $\abs{E(G_0)}$ is the edge count of $G_0$.
\end{lemma}

\begin{lemma}[Telescoping product identity]
\lean{telescope_prod}
\label{telescope_prod}
\leanok
\difficulty{5}
For every integer $n \ge 2$, the telescoping product
\[
  \prod_{i=0}^{n-3} \frac{n - i - 2}{n - i} = \frac{2}{n(n-1)}
\]
holds, where the factors are evaluated as real numbers.
\end{lemma}

\begin{theorem}[Lower bound for the min-cut survival product]
\lean{karger_survival_prob_algorithmic}
\label{karger_survival_prob_algorithmic}
\leanok
\uses{IsMulMinCut, MulCut, Multigraph, Multigraph.vertexCount, telescope_prod}
\difficulty{1}
Let $C$ be a minimum cut of a multigraph $G$, and write $n$ for its number of vertices;
suppose $n \ge 4$. Then
\[
  \frac{2}{n(n-1)} \;\le\; \prod_{i=0}^{n-3} \frac{n-i-2}{n-i}.
\]
\end{theorem}

\begin{theorem}[Karger's contraction algorithm succeeds with high probability]
\lean{karger_whp}
\label{karger_whp}
\leanok
\uses{Multigraph, Multigraph.vertexCount}
\difficulty{5}
Let $G$ be a loopless multigraph whose vertex set is finite, and suppose $G$ has $n \ge
4$ vertices. Then running Karger's random-contraction algorithm $\binom{n}{2} =
n(n-1)/2$ times independently and returning the smallest cut found succeeds with
probability
\[
  1 - \left(1 - \frac{2}{n(n-1)}\right)^{n(n-1)/2} \;\ge\; 1 - e^{-1}.
\]
\end{theorem}

\begin{lemma}[A cardinality bound from a probability sum]
\lean{card_le_of_disjoint_prob_lb}
\label{card_le_of_disjoint_prob_lb}
\leanok
\difficulty{2}
Let $m$ be a natural number and $p > 0$ a real number. If $m p \le 1$, then $m \le
p^{-1}$.
\end{lemma}

\begin{lemma}[Closed form for $\binom{n}{2}$]
\lean{choose_two}
\label{choose_two}
\leanok
\difficulty{3}
For every natural number $n \ge 2$, the binomial coefficient counting the unordered
pairs from an $n$-element set satisfies
\[
  \binom{n}{2} = \frac{n(n-1)}{2},
\]
an identity of real numbers.
\end{lemma}

\begin{lemma}[Power bound on a binomial coefficient]
\lean{choose_le_pow}
\label{choose_le_pow}
\leanok
\difficulty{1}
For all natural numbers $n$ and $\alpha$, the binomial coefficient of $n$ over $2\alpha$
is bounded by the corresponding power of $n$: \[\binom{n}{2\alpha} \le n^{2\alpha}.\]
\end{lemma}

\begin{theorem}[A counting bound of $\binom{n}{2}$]
\lean{num_mincuts_le_choose}
\label{num_mincuts_le_choose}
\leanok
\uses{Multigraph, Multigraph.vertexCount, card_le_of_disjoint_prob_lb, choose_two}
\difficulty{4}
Let $G$ be a multigraph whose (finite) vertex set has $n \ge 2$ elements, and let $m$ be
a natural number. If
\[
  m \cdot \frac{2}{n(n-1)} \le 1,
\]
then $m \le \binom{n}{2}$.
\end{theorem}

\begin{theorem}[At most $n^{2\alpha}$ near-minimum cuts]
\lean{num_alpha_mincuts_le}
\label{num_alpha_mincuts_le}
\leanok
\uses{Multigraph, Multigraph.vertexCount, card_le_of_disjoint_prob_lb, choose_le_pow}
\difficulty{3}
Let $G$ be a multigraph with $n$ vertices, and let $\alpha$ and $m$ be natural numbers
with $\alpha \ge 1$ and $2\alpha \le n$. If, as an inequality of real numbers, $m \cdot
\binom{n}{2\alpha}^{-1} \le 1$, then $m \le n^{2\alpha}$.
\end{theorem}
